\title{DeepSeekMath: Pushing the Limits of Mathematical Reasoning in Open Language Models}

\begin{abstract}
Mathematical reasoning remains a major hurdle for language models due to its inherently structured and intricate demands. In this study, we present DeepSeekMath~7B, developed by further pre-training the DeepSeek-Coder-Base-v1.5 7B model using 120B tokens rich in mathematical content extracted from Common Crawl, as well as integrated natural language and programming data. DeepSeekMath~7B attains a notable 51.7\% score on the MATH benchmark, designed for competition-level mathematics, without the assistance of supplementary toolkits or ensemble voting strategies—narrowing the performance gap with advanced models such as Gemini-Ultra and GPT-4. When evaluating self-consistency over 64 generated samples, DeepSeekMath~7B achieves 60.9\% accuracy on MATH. The model’s proficiency in mathematical reasoning is largely driven by two primary innovations: first, the systematic utilization of open-source web datasets via a carefully crafted data selection framework; second, the development of Group Relative Policy Optimization (GRPO), an adaptation of Proximal Policy Optimization (PPO) that not only elevates mathematical reasoning capabilities but also makes PPO more memory-efficient.
\end{abstract}

\section{Introduction}

Large language models (LLMs) have transformed how artificial intelligence addresses mathematical reasoning, leading to notable progress on quantitative reasoning tasks \citep{MATH} and in geometry-related evaluations \citep{trinh2024solving}. Furthermore, these models have become valuable tools in aiding human problem solvers tackle advanced mathematical challenges \citep{tao}. Nevertheless, leading solutions such as GPT-4 \citep{gpt4} and Gemini-Ultra \citep{gemini} remain closed-source, leaving available open models with substantially lower performance.

This work presents DeepSeekMath, a specialized language model tailored for mathematical tasks, which demonstrates substantially better mathematical proficiency than other open-source models and performs close to GPT-4 on standardized academic evaluations. To facilitate this, we assemble the DeepSeekMath Corpus, an extensive, high-quality pre-training dataset containing 120B math-specific tokens. The corpus is curated from Common Crawl (CC) via a classifier built on fastText \citep{joulin2016fasttext}. Initially, this classifier is trained with positive samples drawn from OpenWebMath \citep{openwebmath} and a varied set of unrelated web content as negative examples. Once trained, the classifier identifies further mathematical instances within the CC data, which are subsequently validated and refined by human annotators. The improved dataset is used to retrain the classifier, thus enhancing its accuracy. Empirical analysis demonstrates the high standard of this corpus, as evidenced by our base model, DeepSeekMath-Base 7B, achieving 64.2\% accuracy on GSM8K \citep{gsm8k} and 36.2\% on the MATH competition dataset \citep{MATH}, surpassing the performance of Minerva 540B \citep{minerva}. Additionally, since the DeepSeekMath Corpus incorporates multiple languages, we observe progress on Chinese mathematics benchmarks \citep{wei2023cmath,agieval}. We suggest that our methodological advancements in curating mathematical data provide a valuable foundation for further investigations and anticipate continued advancements in this area.

DeepSeekMath-Base is built upon DeepSeek-Coder-Base-v1.5 7B \citep{deepseek-coder}, as our experiments indicate that initializing from a model pretrained on code is more advantageous than beginning with a standard language model. Notably, we find that mathematical training not only strengthens mathematical competence but also leads to improved results on general benchmarks such as MMLU \citep{mmlu} and BBH \citep{bbh}, suggesting enhancements in both mathematical and broader reasoning abilities.

Following the pre-training stage, we further refine DeepSeekMath-Base using mathematical instruction tuning. This involves incorporating data reflecting chain-of-thought \citep{cot}, program-of-thought \citep{pot,pal}, and reasoning integrated with tools \citep{tora}. The final product, DeepSeekMath-Instruct 7B, surpasses all comparable 7B models and rivals the performance of state-of-the-art 70B instruction-tuned open-source models.

Additionally, we present Group Relative Policy Optimization (GRPO), a novel reinforcement learning approach inspired by Proximal Policy Optimization (PPO) \citep{schulman2017proximal}. GRPO eliminates the need for a critic model, instead leveraging group-based score estimation to serve as a baseline, thereby markedly lowering resource requirements for training. By restricting the learning process to a selection of English instruction tuning data, GRPO achieves notable gains over the robust DeepSeekMath-Instruct model on both domain-specific (GSM8K: 82.9\% $\rightarrow$ 88.2\%, MATH: 46.8\% $\rightarrow$ 51.7\%) and external mathematical datasets (such as CMATH: 84.6\% $\rightarrow$ 88.8\%) during reinforcement learning.

We also propose a unified framework to conceptualize various learning strategies—including Rejection Sampling Fine-Tuning (RFT) \citep{yuan2023scaling}, Direct Preference Optimization (DPO) \citep{dpo}, PPO, and GRPO—demonstrating that these can be categorized as either explicit or streamlined RL methods. Our systematic empirical analysis spans facets such as online versus offline optimization, process versus outcome supervision, and comparisons between single-turn and iterative RL setups. This thorough investigation elucidates the core constituents of this paradigm. Finally, we discuss the mechanisms by which our RL strategy enhances the effectiveness of instruction-tuned models and outline promising avenues for future research aiming at more efficient RL within this consolidated framework.

\subsection{Contributions}
We make advancements in large-scale mathematical pre-training and conduct an extensive study of reinforcement learning strategies.

\textbf{Math Pre-Training at Scale}

\begin{itemize}[topsep=0pt]
    \item Our investigation reveals that Common Crawl, despite being a general web resource, encompasses substantial mathematically relevant content. Through the design of a rigorous filtering pipeline, we curate the DeepSeekMath Corpus, comprising 120B tokens sourced from web pages with rich mathematical content. This corpus vastly exceeds the scope of prior datasets: it is nearly 7-fold larger than that used by Minerva \citep{minerva} and almost 9 times larger than OpenWebMath \citep{openwebmath}.

    \item DeepSeekMath-Base 7B, our pre-trained model, demonstrates performance on par with Minerva 540B \citep{minerva}, emphasizing that raw parameter count is not the sole determinant of mathematical reasoning proficiency. Robust results can be obtained even with a more compact architecture if exposed to high-quality data during pre-training.

    \item We document insights from our math-focused training regimes, notably that conducting code pre-training before math pre-training substantially enhances model performance on mathematical tasks, both in tool-augmented and standard settings. This finding contributes to the ongoing debate regarding the benefits of code pre-training for reasoning skills: we contend that it is indeed beneficial, particularly for mathematical reasoning.

    \item Contrary to common practice, our results indicate that supplementing pre-training with arXiv papers does not yield discernible gains on the suite of mathematical benchmarks considered herein, despite its frequent adoption in math-oriented research.
\end{itemize}

\textbf{Exploration and Analysis of Reinforcement Learning}

\begin{itemize}[topsep=0pt]
    \item We present Group Relative Policy Optimization (GRPO), a novel reinforcement learning approach that forgoes the traditional critic model. Instead, GRPO leverages group scores to estimate the baseline, thereby markedly reducing the computational and training resources required, especially in comparison to Proximal Policy Optimization (PPO).
    \item Our experiments illustrate that GRPO considerably boosts the performance of our instruction-tuned DeepSeekMath-Instruct model using only instruction-tuning datasets. Additionally, we document measurable gains in generalization to out-of-domain scenarios over the course of the reinforcement learning training.
    \item We outline a comprehensive framework that conceptualizes and relates a range of reinforcement learning methodologies, including RFT, DPO, PPO, and GRPO. Moreover, we undertake thorough empirical studies—including comparisons between online and offline learning, outcome versus process-level supervision, and single-turn as opposed to iterative reinforcement learning—to analyze the fundamental building blocks underpinning this unified paradigm.
    \item Drawing on this integrated perspective, we investigate the underlying mechanisms driving the effectiveness of reinforcement learning, and identify several promising directions for enhancing reinforcement learning outcomes in large language models (LLMs).
\end{itemize}

\subsection{Summary of Evaluations and Metrics}

\begin{itemize}[topsep=0pt]
    \item \textbf{English and Chinese Mathematical Reasoning}: We systematically evaluate our models on a variety of English and Chinese mathematical reasoning datasets, encompassing problem sets from elementary through collegiate levels. For English, these include GSM8K \citep{gsm8k}, MATH \citep{MATH}, SAT \citep{llemma}, OCW Courses \citep{minerva}, and MMLU-STEM \citep{mmlu}; for Chinese, we utilize MGSM-zh \citep{mgsm}, CMATH \citep{wei2023cmath}, Gaokao-MathCloze \citep{agieval}, and Gaokao-MathQA \citep{agieval}. Our evaluation addresses both the models’ capacity to generate comprehensive written solutions unaided by external tools, as well as their proficiency in leveraging Python for problem solving.

    On English benchmarks, DeepSeekMath-Base demonstrates strong performance, achieving results comparable to the proprietary Minerva 540B \citep{minerva}, and significantly outperforming all open-source base models (such as Mistral 7B \citep{mistral} and Llemma-34B \citep{llemma}), regardless of the presence or absence of math-specific pre-training, frequently by substantial margins. Importantly, DeepSeekMath-Base excels in the Chinese benchmark tasks, which can be attributed to our deviation from previous approaches \citep{minerva,llemma} that utilized only English-language pre-training, and our inclusion of diverse, high-quality multilingual data. Upon further instruction tuning and application of reinforcement learning, the enhanced DeepSeekMath-Instruct and DeepSeekMath-RL models achieve state-of-the-art open-source accuracy on challenging competitions such as the MATH dataset, crossing the 50\% threshold for the first time in the open-source landscape.

\item \textbf{Formal Mathematics}: For the informal-to-formal theorem proving benchmark from \citep{dsp_proof}, we assess DeepSeekMath-Base on miniF2F \citep{minif2f} with Isabelle \citep{isabelle} serving as the proof assistant. Our findings indicate that DeepSeekMath-Base achieves robust performance in few-shot autoformalization.

\item \textbf{Natural Language Understanding, Reasoning, and Code}: To thoroughly characterize the general capabilities of the models in comprehension, reasoning, and coding, DeepSeekMath-Base is benchmarked on the Massive Multitask Language Understanding (MMLU) dataset \citep{mmlu}, which features 57 diverse multiple-choice categories, as well as on BIG-Bench Hard (BBH) \citep{bbh}, which contains 23 intricate tasks that frequently demand multi-step reasoning. Furthermore, HumanEval \citep{codex} and MBPP \citep{mbpp} are included to benchmark code generation skills. The outcomes reveal that mathematical pre-training bolsters both general language understanding and reasoning proficiency.

\section{Math Pre-Training}

\subsection{Data Collection and Decontamination} This section details our approach for developing the DeepSeekMath Corpus utilizing Common Crawl data. Illustrated in Figure \ref{fig:data_collect}, we describe an iterative framework that enables the systematic compilation of a large mathematical corpus from Common Crawl, initiated by a seed dataset consisting of a limited yet highly curated set of mathematics-focused data. It should be emphasized that this methodology can extend to other domains, such as software engineering.

To initiate this process, we utilize OpenWebMath \citep{openwebmath}, a repository of curated mathematical web texts, as the foundation for our seed corpus. Leveraging this, we train a fastText classifier \citep{joulin2016fasttext} to retrieve additional mathematical web content reminiscent of OpenWebMath. Concretely, we randomly sample 500,000 entries from the seed dataset as positive examples and supplement these with 500,000 negative examples drawn from Common Crawl. Training is conducted using the public fastText library\footnote{\url{https://fasttext.cc}} with hyperparameters: vector dimension set at 256, learning rate of 0.1, maximum word n-gram length of 3, minimum word frequency of 3, and 3 epochs. To further pare down Common Crawl, we apply both URL-based deduplication and near-duplicate filtering, narrowing it to 40B HTML web pages. Employing the trained fastText model, we then recall likely mathematical web content from this reduced dataset. To ensure quality, collected pages are ranked according to their fastText score, and only the top-scoring subset is retained. The optimal dataset size is determined via pre-training trials across datasets containing the top 40B, 80B, 120B, and 160B tokens, with the first cycle retaining the top 40B tokens.

Following the initial round of data gathering, a significant portion of mathematical web pages remains uncollected, primarily because the fastText model is initially trained on a collection of positive samples that lacks adequate variability. To address this, we expand the seed corpus by incorporating further mathematical web domains, aiming to enhance the effectiveness of the fastText model. Concretely, we segment the full Common Crawl dataset into separate domains, where each domain includes web pages with an identical base URL. For each domain, we determine the proportion of web pages acquired during the first round of collection. Domains where more than 10\% of pages are captured are classified as math-centric (such as \url{mathoverflow.net}). Next, we conduct a manual review of URLs in these math-centric domains to pinpoint those directly tied to mathematical material (for example, \url{mathoverflow.net/questions}). Web pages affiliated with these URLs that have not yet been collected are added to the seed set. This strategy provides a richer array of positive samples for the fastText model, improving its ability to identify and retrieve additional mathematical pages in subsequent cycles. Across four data collection cycles, this process yields a total of 35.5M mathematical web pages, amassing 120B tokens. By the fourth cycle, we find that approximately 98\% of the material was already acquired in the preceding round, leading us to conclude the data collection at that point.

To prevent overlap with evaluation datasets, we apply the methodology described in \cite{deepseek-coder} to exclude any web pages containing content from well-known English mathematical benchmarks such as GSM8K \citep{gsm8k} and MATH \citep{MATH}, as well as Chinese benchmarks such as CMATH \citep{wei2023cmath} and AGIEval \citep{agieval}. The exclusion rule is as follows: any passage that contains a 10-gram substring perfectly matching a segment from the benchmark datasets is removed from our training set. For benchmark passages comprising fewer than 10 but at least 3 grams, we remove the corresponding pages if they match exactly.

\subsection{Validating the Quality of the DeepSeekMath~Corpus}

We conduct pre-training experiments to evaluate the DeepSeekMath~Corpus against other recent corpora developed for mathematical machine learning:

\begin{itemize}[topsep=0pt]
    \item \textbf{MathPile} \citep{mathpile}: a composite dataset (8.9B tokens) assembled from diverse sources, including textbooks, Wikipedia, ProofWiki, CommonCrawl, StackExchange, and arXiv, with arXiv comprising the largest share (over 85\%);
    \item \textbf{OpenWebMath} \citep{openwebmath}: a subset of CommonCrawl curated to focus on mathematical web content, amounting to 13.6B tokens;
    \item \textbf{Proof-Pile-2} \citep{llemma}: an amalgamated mathematical corpus that combines OpenWebMath, AlgebraicStack (containing 10.3B tokens of mathematical code), and arXiv manuscripts (28.0B tokens). In experiments involving Proof-Pile-2, we adhere to the arXiv:Web:Code 2:4:1 proportioning outlined by \cite{llemma}.
\end{itemize}

\subsubsection{Training Setting}
\label{sec:quality-policy}
We specialize a 1.3B-parameter general-purpose pre-trained language model—built on the same architecture as the DeepSeek LLMs \citep{deepseek-llm}, herein referred to as DeepSeek-LLM 1.3B—by further training it with mathematical data. Each mathematical dataset is used to fine-tune an individual model for a total of 150B tokens. The training procedure utilizes the resource-efficient HAI-LLM framework \citep{haillm}. Adhering to the methodologies established in the DeepSeek LLMs, we adopt the AdamW optimizer \citep{adamW} with hyperparameters $\beta_1=0.9$, $\beta_2=0.95$, and $\mathrm{weight\_decay}=0.1$. The learning rate follows a staged schedule: it is linearly warmed up over the initial 2,000 steps to its maximum, reduced to 31.6\% of this value after 80\% of the training period, and finally annealed to 10.0\% of the peak learning rate after 90\%. The peak learning rate is capped at 5.3e-4, the batch size is set to 4M tokens, and models are trained using a context window of 4K tokens.

\subsubsection{Evaluation Results}

\textbf{The DeepSeekMath~Corpus stands out in terms of data quality, multilingual breadth, and corpus size.}

\begin{itemize}[topsep=0pt]
    \item \textbf{High-quality}:
    We assess the downstream abilities of the model on eight math-focused evaluation benchmarks utilizing few-shot chain-of-thought prompting \cite{cot}. According to Table \ref{tab:corpora_comparison}, the model fine-tuned on the DeepSeekMath~Corpus achieves the best results. Further, Figure \ref{fig:corpora_comparisons} reveals that a model trained on DeepSeekMath outperforms a model trained on Proof-Pile-2 even after 50B tokens (equivalent to one complete epoch over Proof-Pile-2), suggesting that the overall quality of DeepSeekMath is superior. 
    \item \textbf{Multilingual}:
    DeepSeekMath~Corpus contains mathematical texts in multiple languages, with English and Chinese constituting the largest language portions. Table \ref{tab:corpora_comparison} demonstrates that leveraging the DeepSeekMath~Corpus significantly bolsters mathematical reasoning proficiency in both English and Chinese. By comparison, previous mathematical corpora, which are predominantly monolingual (English), provide minimal benefit—or can even adversely affect—performance on mathematical reasoning tasks in Chinese.
    \item \textbf{Large-scale}:
    DeepSeekMath~Corpus exceeds other mathematical datasets by several multiples in size. As presented in Figure \ref{fig:corpora_comparisons}, DeepSeek-LLM 1.3B, once trained with DeepSeekMath, exhibits more rapid and sustained improvements on evaluation benchmarks. In contrast, models trained on baseline corpora, which are considerably smaller and exhausted through several epochs within training, see their performance gains quickly saturate.
\end{itemize}

\subsection{Training and Evaluating DeepSeekMath-Base 7B}

This section presents DeepSeekMath-Base 7B, a foundational model designed with robust mathematical reasoning capacities. The model initialization leverages DeepSeek-Coder-Base-v1.5 7B \citep{deepseek-coder}, followed by pretraining over 500 billion tokens. The training dataset is composed as follows: 56\% DeepSeekMath Corpus, 4\% AlgebraicStack, 10\% arXiv, 20\% Github code, and 10\% natural language texts from Common Crawl (in both English and Chinese). The training configuration closely follows the procedure detailed in Section \ref{sec:quality-policy}, with two key deviations: the learning rate is set with a peak value of 4.2e-4 and the batch size comprises 10 million tokens.

A thorough evaluation is performed to measure the mathematical proficiency of DeepSeekMath-Base 7B. Assessment dimensions include its ability to generate self-sufficient mathematical answers independently, utilize external tools in problem solving, and execute formal theorem proofs. In addition to its mathematical focus, we also present an overview of the model's broader capabilities—spanning natural language understanding, logical reasoning, and programming aptitude.

\paragraph{Mathematical Problem Solving with Step-by-Step Reasoning}

DeepSeekMath-Base 7B is systematically tested on its capacity for stepwise mathematical problem solving using few-shot chain-of-thought prompting \citep{cot}. The evaluation spans eight benchmarks in both English and Chinese, incorporating domains such as quantitative reasoning (for instance, GSM8K \citep{gsm8k}, MATH \citep{MATH}, and CMATH \citep{wei2023cmath}) as well as multiple-choice assessments (such as MMLU-STEM \citep{mmlu} and Gaokao-MathQA \citep{agieval}). These benchmarks represent a wide spectrum of mathematical disciplines, ranging from foundational to university-level difficulty.

Table \ref{tab:base_math_cot} highlights that DeepSeekMath-Base 7B consistently achieves the highest results among all eight evaluated benchmarks when compared with other open-source base models. This includes not only the extensively adopted generalist Mistral 7B \citep{mistral} but also Llemma 34B \citep{llemma}, a recent model specializing in mathematical reasoning and trained with Proof-Pile-2 \citep{llemma}. A particularly notable result is observed on the challenging MATH dataset, where DeepSeekMath-Base attains more than a 10\% absolute improvement over all open-source counterparts. Furthermore, it exceeds the performance of Minerva 540B \citep{minerva}—a proprietary model derived from PaLM \citep{palm} and trained on math-specific corpora—despite Minerva’s considerably larger scale.

\paragraph{Mathematical Problem Solving with Tool Use} To assess models’ capacity for program-assisted mathematical reasoning, we conduct few-shot program-of-thought prompting \citep{pot,pal} on GSM8K and MATH datasets. In this protocol, each task requires the model to author a Python script, leveraging libraries such as \textit{math} and \textit{sympy} to manage complex calculations. The program’s output is adopted as the final answer. According to Table \ref{tab:base_math_pot_proof}, DeepSeekMath-Base 7B surpasses the previous top-performing Llemma 34B in these evaluations.

\paragraph{Formal Mathematics}

The automation of formal proofs is instrumental for verifying the precision and correctness of mathematical arguments, while also streamlining proof development—a topic receiving increasing focus. We test DeepSeekMath-Base 7B using the informal-to-formal proof conversion task described in \citep{dsp_proof}, which involves transforming an informal theorem, its formal statement, and its informal proof into a formalized proof. Evaluation is carried out on the miniF2F \citep{minif2f} dataset, which targets formal mathematical problems at Olympiad level, by generating formal Isabelle proofs using few-shot prompting. Consistent with the approach in \cite{dsp_proof}, we prompt the model to provide proof sketches and employ the Sledgehammer tool \citep{sledgehammer} to complete any missing justification steps. Table \ref{tab:base_math_pot_proof} demonstrates that DeepSeekMath-Base 7B achieves strong results in automated formal proof generation.

\paragraph{Natural Language Understanding, Reasoning, and Code}

We evaluate the model’s comprehension using MMLU \citep{mmlu}, logical reasoning via BBH \citep{bbh}, and programming capabilities with the HumanEval \citep{codex} and MBPP \citep{mbpp} benchmarks. Table \ref{tab:understanding_reasoning_code} reveals that DeepSeekMath-Base 7B considerably advances performance on MMLU and BBH compared with its predecessor, DeepSeek-Coder-Base-v1.5 \citep{deepseek-coder}, underscoring the advantageous effect of mathematical pretraining on language and reasoning abilities. Moreover, thanks to continual training involving code data, DeepSeekMath-Base 7B preserves the programming proficiency levels of DeepSeek-Coder-Base-v1.5 across both coding benchmarks. Collectively, DeepSeekMath-Base 7B outstrips the generalist Mistral 7B \citep{mistral} across all three evaluated domains: reasoning and code generation.

\section{Supervised Fine-Tuning}

\subsection{SFT Data Curation}

We assembled a dataset for mathematical instruction tuning that encompasses both English and Chinese mathematical tasks, sampling from a wide array of mathematical disciplines and a spectrum of problem difficulties. Each problem instance is paired with detailed solutions using chain-of-thought (CoT) \citep{cot}, program-of-thought (PoT) \citep{pot,pal}, and tool-integrated reasoning approaches \citep{tora}. The curated dataset comprises 776,000 training samples.

\begin{itemize}[topsep=0pt]
    \item \textbf{English mathematical datasets}: Tool-integrated solution annotations were added to GSM8K and MATH datasets. Additionally, we incorporated a filtered subset of MathInstruct \citep{MathInstruct} and included the training portion of Lila-OOD \citep{lila}, featuring CoT and PoT solutions. Our English dataset selection spans multiple mathematical disciplines, including but not limited to algebra, probability, number theory, calculus, and geometry.
    \item \textbf{Chinese mathematical datasets}: For Chinese K-12 mathematics, we gathered problems from 76 different sub-domains (e.g., linear equations) and supplied answers formulated both as CoT explanations and in a tool-integrated reasoning style.
\end{itemize}

\subsection{Training and Evaluating DeepSeekMath-Instruct 7B}

This section details the supervised instruction tuning of DeepSeekMath-Instruct 7B, built upon DeepSeekMath-Base. During training, we concatenated random samples up to a maximum sequence length of 4,000 tokens. The model was trained for 500 iterations with a mini-batch size of 256 and a fixed learning rate set to 5e-5.

We conducted an assessment of the model’s quantitative reasoning performance, evaluating both its unaided abilities and its capabilities with tool support, across four English and Chinese benchmarks. To contextualize our results, we compare performance against several leading contemporary models:

\begin{itemize}[topsep=0pt]
    \item \textbf{Closed-source models} surveyed include: (1) various iterations from the GPT family, with GPT-4 \citep{gpt4} and GPT-4 Code Interpreter~\footnote{\url{https://openai.com/blog/chatgpt-plugins\#\#code-interpreter}} recognized as particularly strong, (2) Gemini Ultra and Gemini Pro \citep{gemini}, (3) Inflection-2 \citep{inflection-2}, (4) Grok-1~\footnote{\url{https://x.ai/model-card}}, and additional models by Chinese enterprises, namely (5) Baichuan-3~\footnote{\url{https://www.baichuan-ai.com}}, (6) the newly released GLM-4~\footnote{\url{https://open.bigmodel.cn/dev/api\#glm-4}} 

    which are part of the GLM family \citep{glm}.
\end{itemize}

The majority of these models are designed for general use and have been subjected to various alignment processes.

\item \textbf{Open-source models} feature a mix of generalist and mathematics-augmented models: (1) DeepSeek-LLM-Chat 67B \citep{deepseek-llm}, (2) Qwen 72B \citep{qwen}, (3) SeaLLM-v2 7B \citep{seallm}, and (4) ChatGLM3 6B \citep{chatglm3} as general-purpose systems, and additional models with mathematical enhancements such as (5) InternLM2-Math 20B~\footnote{\url{https://github.com/InternLM/InternLM-Math}}, an InternLM2 derivative trained specifically on mathematics and subsequently instruction-tuned, (6) Math-Shepherd-Mistral 7B, which leverages PPO training \citep{schulman2017proximal} on Mistral 7B \citep{mistral} using a process-level supervised reward model, (7) the WizardMath family \citep{wizardmath}, which boosts mathematical capabilities of Mistral 7B and Llama-2 70B \citep{llama2} through evolve-instruct (i.e., AI-generated instruction tuning) and PPO, primarily drawing training data from GSM8K and MATH, (8) MetaMath 70B \citep{metamath}, a Llama-2 70B fine-tuned on GSM8K and MATH with data augmentation, (9) ToRA 34B \cite{tora}, a CodeLlama 34B variant refined for tool-supported mathematical reasoning, and (10) MAmmoTH 70B \citep{MathInstruct}, which applies MathInstruct-based instruction tuning to Llama-2 70B.
\end{itemize}

According to Table \ref{tab:sft_rl_math}, when tool-assisted methods are not permitted in evaluation, DeepSeekMath-Instruct 7B excels at stepwise reasoning. Particularly on the challenging MATH dataset, it outperforms every open-source model and most proprietary systems (such as Inflection-2 and Gemini Pro) by a margin of at least 9 percentage points. This superiority holds even in comparison to significantly larger models like Qwen 72B or those optimized with mathematics-specific reinforcement learning, such as WizardMath-v1.1 7B. DeepSeekMath-Instruct demonstrates results on par with Chinese closed-source models like GLM-4 and Baichuan-3 in MATH, though it remains behind leading models such as GPT-4 and Gemini Ultra.

In settings where integration of both natural language reasoning and code-based tools for problem-solving is permitted, DeepSeekMath-Instruct 7B reaches an accuracy near 60\% on the MATH benchmark, surpassing all other available open-source models. On additional benchmarks, this model stands competitive with DeepSeek-LLM-Chat 67B, the previously leading open-source model, despite being an order of magnitude smaller.

\section{Reinforcement Learning}

\subsection{Group Relative Policy Optimization} Recent research has shown that reinforcement learning (RL) can further enhance the mathematical reasoning capabilities of LLMs beyond what is achievable through Supervised Fine-Tuning (SFT) alone \citep{wang2023math,wizardmath}. This section presents our novel and resource-efficient RL approach, termed Group Relative Policy Optimization (GRPO).

\subsubsection{From PPO to GRPO} Proximal Policy Optimization (PPO) \citep{schulman2017proximal} serves as a widely adopted actor-critic RL framework for tuning LLMs in the RL phase \citep{ouyang2022training}. PPO updates models by optimizing a surrogate objective, formulated as:

\begin{equation}
\footnotesize
    \mathcal{J}_{PPO}(\theta) = \mathbb{E}{[q \sim P(Q), o \sim \pi_{\theta_{old}}(O|q)]} \frac{1}{|o|} \sum_{t=1}^{|o|} \min \left[ \frac{\pi_\theta(o_{t} | q, o_{<t})}{\pi_{\theta_{old}}(o_{t} | q, o_{<t})} A_{t}, \text{clip} \left( \frac{\pi_\theta(o_{t} | q, o_{<t})}{\pi_{\theta_{old}}(o_{t} | q, o_{<t})}, 1 - \varepsilon, 1 + \varepsilon \right)  A_{

Here, $\pi_{\theta}$ and $\pi_{\theta_{old}}$ denote the policy model being trained and its previous version, respectively. The variables $q$ and $o$ refer to questions and sampled outputs, with questions drawn from the question dataset and outputs generated by the old policy $\pi_{\theta_{old}}$. The parameter $\varepsilon$ serves as a clipping hyper-parameter within PPO to ensure stability in training. The term $A_t$ stands for the advantage, which is estimated using Generalized Advantage Estimation (GAE) \citep{gae}, computed with reward signals $\{r_{\ge t}\}$ and the learned value function $V_{\psi}$. As a result, PPO jointly trains a value function alongside the policy model, and to prevent the reward model from being over-optimized, it is standard to introduce a KL divergence penalty per token between the policy and a reference model within the reward, following the approach of \citep{ouyang2022training}:

\begin{equation}
    r_{t} = r_\varphi(q, o_{\le t}) - \beta \log\frac{\pi_{\theta}(o_{t}|q, o_{<t})}{\pi_{ref}(o_{t}|q, o_{<t})},
\label{eq:PPO-reward}
\end{equation}

where $r_\varphi$ represents the reward model, $\pi_{ref}$ designates the reference policy (often the original SFT model), and $\beta$ determines the weight of the KL regularization.

Maintaining a separate value function of a similar scale as the policy model, as required by PPO, significantly increases computational and memory demands. Furthermore, in RL training, the value function provides a baseline for advantage calculations to reduce variance. However, in LLM applications, the reward model generally provides scores only for the final token, complicating the training of a token-level accurate value function. To overcome these limitations, as depicted in Figure \ref{fig:grpo}, we introduce Group Relative Policy Optimization (GRPO). GRPO eliminates the need for an explicit value function approximation as in PPO by leveraging the mean reward from several outputs sampled in response to the same prompt as a baseline. Specifically, for each question $q$, GRPO samples a collection $\{o_1, o_2, \cdots, o_G\}$ from $\pi_{\theta_{old}}$, and optimizes the policy through the objective below:

\begin{equation}
\footnotesize
\begin{split}
    \mathcal{J}_{GRPO}(\theta) &= \mathbb{E}{[q \sim P(Q), \{o_i\}_{i=1}^G \sim \pi_{\theta_{old}}(O|q)]}  \\
    & \frac{1}{G}\sum_{i=1}^G\frac{1}{|o_i|} \sum_{t=1}^{|o_i|} \left\{ \min \left[ \frac{\pi_\theta(o_{i,t} | q, o_{i,<t})}{\pi_{\theta_{old}}(o_{i,t} | q, o_{i,<t})} \hat{A}_{i,t}, \text{clip} \left( \frac{\pi_\theta(o_{i,t} | q, o_{i,<t})}{\pi_{\theta_{old}}(o_{i,t} | q, o_{i,<t})}, 1 - \varepsilon, 1 + \varepsilon \right)  \hat{A}_{i,t} \right] - \beta \mathbb{D}_{KL}\left[\pi_{\theta} || \pi_{ref}\right]\right\} ,
\end{split}
\label{eq:GRPO-obj}
\end{equation}

with $\varepsilon$ and $\beta$ serving as tunable hyper-parameters, and $\hat{A}_{i,t}$ denoting the advantage, computed based on relative rewards among members of each group—further details are given in the subsequent subsections. By adopting this group-relative approach to advantage computation, GRPO naturally mirrors the reward models’ structure, which are often trained using pairwise output comparisons for the same question. Additionally, GRPO imposes the KL divergence between the trained and reference policies directly onto the objective function, instead of integrating the KL term into the reward. This circumvents the need for incorporating the penalty into the advantage computation, thereby simplifying $\hat{A}_{i,t}$'s calculation. Contrary to the KL penalty in (\ref{eq:PPO-reward}), GRPO estimates the KL divergence using the following unbiased estimator, as described in \citep{kl_approx

\begin{algorithm}[t]
  \small
  \caption{Iterative Group Relative Policy Optimization}
  \textbf{Input} initial policy model $\pi_{\theta_{\text{init}}}$; reward models $r_\varphi$; task prompts $\mathcal{D}$; 
  hyperparameters $\varepsilon$, $\beta$, $\mu$
  \begin{algorithmic}[1]
    \State Initialize policy model $\pi_\theta$ with $\pi_{\theta_{\text{init}}}$
    \For{iteration = 1, \dots, I}
       \State Set reference model $\pi_{ref}$ as $\pi_{\theta}$
      \For{step = 1, \dots, M}
      \State Draw a minibatch $\mathcal{D}_b$ from $\mathcal{D}$
      \State Assign $\pi_{\theta_{old}} \leftarrow \pi_{\theta}$ as the prior policy 
      \State For each question $q \in \mathcal{D}_b$, generate $G$ candidate responses $\{o_i\}_{i=1}^G \sim \pi_{\theta_{old}} (\cdot \mid q) $
      \State Evaluate rewards $\{r_i\}_{i=1}^{G}$ for each output $o_i$ via $r_{\varphi}$
      \State Compute group relative advantage $\hat{A}_{i,t}$ for the $t$-th token of $o_i$
      \For{GRPO iteration = 1, \dots, $\mu$}
        \State Adjust $\pi_{\theta}$ by maximizing the GRPO loss function (see Equation \ref{eq:GC-GRPO})
      \EndFor
    \EndFor
    \State Periodically update $r_\varphi$ with continuous learning through replay buffer techniques.
    \EndFor
  \end{algorithmic}
  \textbf{Output} $\pi_\theta$
  \label{alg:iter-grpo}
\end{algorithm}

\subsubsection{Outcome Supervision RL with GRPO} More precisely, for each question $q$, we sample a collection of candidate outputs $\{o_1, o_2, \cdots, o_G\}$ from the previous policy $\pi_{\theta_{old}}$. Using the reward model, each candidate is assigned a reward, resulting in a set $\mathbf{r}=\{r_1, r_2, \cdots, r_G\}$. The raw rewards are then standardized: each reward is adjusted by subtracting the batch mean and dividing by the batch standard deviation. Outcome supervision attributes the resulting normalized reward to the final token of $o_i$ and assigns the same value $\hat{A}_{i, t}$ to all tokens in $o_i$, that is, $\hat{A}_{i, t} = \widetilde{r}_i = \frac{r_i- {\rm mean}(\mathbf{r})}{{\rm std}(\mathbf{r})}$, across the sequence. Optimization of the policy is then performed by maximizing the objective function as detailed in equation (\ref{eq:GRPO-obj}).

\subsubsection{Process Supervision RL with GRPO} Limiting feedback to only the terminal output, as in outcome supervision, may result in insufficient guidance, particularly for intricate mathematical reasoning. In line with \cite{wang2023math}, we therefore examine process supervision, which evaluates each intermediate reasoning stage with dedicated feedback. Given a prompt $q$ and corresponding $G$ samples $\{o_1, o_2, \cdots, o_G\}$, a process-oriented reward model assigns step-level rewards as $\mathbf{R} = \{ \{r_1^{index(1)},\cdots,r_1^{index(K_1)}\}, \cdots,  \{r_G^{index(1)},\cdots,r_G^{index(K_G)}\} \}$, where $index(j)$ marks the last token index of the $j$-th reasoning step and $K_i$ is the number of steps in $o_i$. Each step’s reward is normalized as $\widetilde{r}_i^{index(j)} = \frac{r_i^{index(j)} - {\rm mean(\mathbf{R})}}{{\rm std(\mathbf{R})}}$. To determine the tokenwise advantage, the sum of all subsequent normalized rewards from step indices greater than or equal to the current token index is used: $\hat{A}_{i, t} = \sum_{index(j) \ge t} \widetilde{r}_i^{index(j)}$. The policy is then improved by maximizing the objective in equation (\ref{eq

\subsubsection{Iterative RL with GRPO} As the reinforcement learning process advances, reliance on an outdated reward model may become insufficient for guiding the evolving policy model. To address this, we investigate iterative RL with GRPO. According to Algorithm \ref{alg:iter-grpo}, the iterative GRPO approach involves regenerating training datasets for the reward model based on samples produced by the latest policy model, and updating the reward model continuously using a replay mechanism that retains 10\% of prior data. Subsequently, the reference model is set to match the current policy model, which is then further optimized with the updated reward model.

\subsection{Training and Evaluating DeepSeekMath-RL}

RL is performed utilizing DeepSeekMath-Instruct 7B as the base. For RL training, we utilize chain-of-thought-style problems from GSM8K and MATH extracted from the SFT dataset, amounting to approximately 144K questions. Other SFT instances are omitted to analyze RL's effect on benchmarks lacking direct representation during RL. Reward model training sets are curated in accordance with \citep{wang2023math}. The initial reward model is derived from DeepSeekMath-Base 7B using a learning rate of 2e-5. Within GRPO, we apply a learning rate of 1e-6 to the policy model and set the KL coefficient to 0.04. Each question is sampled for $64$ candidate responses. Maximum sequence length is restricted to 1024, and training batch size is also 1024. Policy model parameters are updated once per exploration phase. For evaluation, DeepSeekMath-RL 7B is assessed on the same benchmarks as DeepSeekMath-Instruct 7B. Here, GSM8K and MATH tasks featuring chain-of-thought reasoning represent in-domain challenges, while all remaining benchmarks are treated as out-of-domain.

Table \ref{tab:sft_rl_math} presents the comparative results of open- and closed-source models employing chain-of-thought and tool-augmented reasoning techniques across English and Chinese benchmarks. Our observations reveal: 1) DeepSeekMath-RL 7B achieves accuracy rates of 88.2\% on GSM8K and 51.7\% on MATH with chain-of-thought reasoning. These results not only lead all open-source models within the 7B to 70B parameter range but also outperform most closed-source counterparts. 2) Importantly, DeepSeekMath-RL 7B is trained exclusively on chain-of-thought style instruction tuning datasets from GSM8K and MATH, with DeepSeekMath-Instruct 7B serving as its initialization. Despite its limited training data domain, DeepSeekMath-RL 7B consistently outperforms DeepSeekMath-Instruct 7B in all measured criteria, thereby highlighting the advantages conferred by reinforcement learning.

\section{Discussion}
This section details insights gleaned from our work in both pre-training and reinforcement learning experiments.

\subsection{Lessons Learnt in Pre-Training}

We begin by outlining observations made during the pre-training process. Unless otherwise indicated, the pre-training configurations adhere to those specified in Section \ref{sec:quality-policy}. In this discussion, references to the DeepSeekMath Corpus pertain to an 89B-token dataset compiled during the project’s second data collection cycle.

\subsubsection{Code Training Benefits Mathematical Reasoning}

There is a widely held, though yet to be empirically confirmed, belief that training models on code enhances their reasoning abilities. Our investigation addresses this hypothesis within the mathematical context, finding that exposure to code indeed improves mathematical reasoning skills—both with tool assistance and without.

To systematically examine the impact of code pre-training on mathematical reasoning performance, we designed experiments using both two-stage and one-stage training paradigms:

\textbf{Two-Stage Training}

\begin{itemize}[topsep=0pt]
    \item \textbf{Code Training for 400B Tokens $\rightarrow$ Math Training for 150B Tokens}: DeepSeek-LLM 1.3B is first trained on 400B tokens consisting of code, followed by an additional 150B tokens focused on mathematical content;
    \item \textbf{General Training for 400B Tokens $\rightarrow$ Math Training for 150B Tokens}: For comparison, we conduct a variant where the initial 400B tokens are drawn from a comprehensive general corpus (compiled by DeepSeek-AI) rather than code, with the intention of evaluating the relative benefits of code versus general corpus pretraining on the development of mathematical reasoning capabilities.
\end{itemize}

\textbf{One-Stage Training}

\begin{itemize}[topsep=0pt]
    \item \textbf{Math Training for 150B Tokens}: DeepSeek-LLM 1.3B is directly trained on 150B tokens drawn from mathematical sources;
    \item \textbf{Training on a mixture of 400B Code Tokens and 150B Math Tokens}: Sequential math training after extensive code training tends to impair coding proficiency. To address this, we examine whether a joint one-stage training regimen—using a mixture of code and math tokens—can still bolster mathematical reasoning and at the same time mitigate catastrophic forgetting.
\end{itemize}

\paragraph{Results}
The empirical outcomes, presented in Table \ref{tab:code-math} and Table \ref{tab:code-math-general-eval}, highlight the impact of these training strategies on downstream evaluation tasks.

Pretraining with code tokens contributes to enhancements in program-aided mathematical reasoning, regardless of whether a two-stage or mixed one-stage training procedure is followed. Table \ref{tab:code-math} indicates that, with two-stage training, merely pretraining on code substantially increases performance on GSM8K and MATH benchmarks involving Python. Applying a subsequent math-focused phase brings further progress. Remarkably, when code and math tokens are combined for joint one-stage training, the detrimental effects of catastrophic forgetting observed with pure two-stage training are alleviated, yielding improvements not only in coding proficiency (Table \ref{tab:code-math-general-eval}) but also in program-aided mathematical problem solving (Table \ref{tab:code-math}).

Moreover, pretraining with code tokens demonstrates a positive impact on mathematical reasoning tasks that do not utilize auxiliary tools. In the two-stage paradigm, starting with code training leads to moderate improvements, and also increases the effectiveness of the later mathematical fine-tuning, achieving the highest observed scores. In contrast, joint one-stage training with both code and math data leads to diminished performance for mathematical reasoning tasks that exclude tool usage. A possible explanation is that DeepSeek-LLM 1.3B, being relatively compact in parameter size, is insufficiently expressive to fully internalize both programming and mathematical knowledge simultaneously.

\subsubsection{ArXiv Papers Show Limited Effectiveness in Enhancing Mathematical Reasoning}

Incorporation of ArXiv papers is a standard practice in assembling math-focused pre-training corpora \citep{minerva,gpt-f,llemma,mathpile}. Despite their prevalence, there remains a lack of thorough investigation into their concrete effects on mathematical reasoning performance. Somewhat unexpectedly, our empirical results indicate that ArXiv papers contribute minimally, if at all, to advancements in mathematical reasoning. We evaluated this by experimenting with multiple model architectures of varying scale, including DeepSeek-LLM 1.3B and DeepSeek-Coder-Base-v1.5 7B \citep{deepseek-coder}, trained on ArXiv datasets subject to distinct preprocessing steps:

\begin{itemize}[topsep=0pt]
    \item \textbf{MathPile} \citep{mathpile}: comprising 8.9B tokens assembled via systematic cleaning and heuristic filtering, of which over 85\% originate from scientific arXiv manuscripts;
    \item \textbf{ArXiv-RedPajama} \citep{redpajama}: consisting of all available arXiv LaTeX submissions with non-content elements—such as preambles, commentary, macro definitions, and bibliographies—excluded, resulting in a dataset of 28.0B tokens.
\end{itemize}

Our study involved individually pre-training DeepSeek-LLM 1.3B with 150B tokens and DeepSeek-Coder-Base-v1.5 7B with 40B tokens on each ArXiv corpus variant. The findings reveal that pre-training solely on ArXiv material fails to deliver measurable gains in mathematical reasoning capability for these models; in certain scenarios, there is observable stagnation or even regression across an array of mathematical evaluation tasks. The assessed tasks span a range of difficulty levels and types, such as quantitative reasoning benchmarks (GSM8K and MATH, referenced in Table \ref{tab:arxiv-cot}), multiple-choice assessments like MMLU-STEM (see Table \ref{tab:arxiv-cot}), and formal proof-oriented tasks exemplified by miniF2F (Table \ref{tab:arxiv_atp}).

It is important to acknowledge that these conclusions are subject to certain caveats. Specifically, the current analysis does not address:

\begin{itemize}[topsep=0pt]
    \item Potential impacts of ArXiv data on other mathematics-related capabilities outside the scope of our benchmarks, for example, informalization of formal mathematical content;
    \item The interplay between ArXiv-sourced data and other modalities or textual sources within a training mix;
    \item Possible advantages of ArXiv datasets that might emerge at substantially larger model scales.
\end{itemize}

Accordingly, more nuanced investigations are warranted, which we defer to subsequent research.

\subsection{Insights into Reinforcement Learning}
\subsubsection{Toward a Unified Paradigm} This section outlines a cohesive analytical framework to interpret diverse training approaches such as SFT, RFT, DPO, PPO, GRPO, and proceeds to empirically examine the core elements of this unified formulation. Broadly, the gradient of a training objective with respect to the parameter $\theta$ in these methodologies can be formalized as:

\begin{equation}
    \nabla_{\theta}\mathcal{J}_{\textcolor{red}{\mathcal{A}}}(\theta) = \mathbb{E}[\underbrace{(q,o) \sim \textcolor{red}{\mathcal{D}}}_{Data \ Source}]\left( \frac{1}{|o|} \sum_{t=1}^{|o|}  \underbrace{GC_{{\mathcal{A}}}(q, o, t, \textcolor{red}{\pi_{{rf}}})}_{Gradient \ Coefficient}  \nabla_{\theta}\log \pi_{\theta}(o_t | q, o_{<t})\right).
\label{eq:objective}
\end{equation}

This framework is comprised of three fundamental elements: 1) the \textit{Data Source} $\mathcal{D}$, which provides the examples for model training; 2) the \textit{Reward Function} $\pi_{{rf}}$, serving as the source for the reward feedback signal used during optimization; and 3) the \textit{Algorithm} $\mathcal{A}$, which translates the available data and associated rewards into a gradient coefficient $GC$, determining the degree of adjustment—either punitive or incentivizing—applied to the learning process. We discuss several canonical strategies within this generalized setting:

\begin{itemize}
    \item  \textbf{Supervised Fine-tuning (SFT)}: In this method, a pretrained model undergoes additional training using curated SFT datasets selected by humans.
    \item \textbf{Rejection Sampling Fine-tuning (RFT)}: RFT builds on the SFT model by further fine-tuning it on a subset of its own generated responses, specifically retaining only those outputs that are correct, based on the SFT prompts.
    \item \textbf{Direct Preference Optimization (DPO)}: DPO further adjusts the SFT model, using a pairwise DPO loss, to optimize on enhanced response sets sampled from the SFT model.
    \item \textbf{Online Rejection Sampling Fine-tuning (Online RFT)}: In contrast to traditional RFT, Online RFT updates the model in real time by initializing from the SFT model and then fine-tuning it continuously on new outputs generated by the current policy.
    \item \textbf{PPO/GRPO}: Here, the model starts from the SFT initialization and is improved through reinforcement with outputs sampled from the active policy at each step.
\end{itemize}

The elements constituting these approaches are outlined in Table \ref{tab:data_policy}. For an expanded explanation of the derivation, please consult Appendix \ref{app:analysis-rl}.

\paragraph{Remarks on Data Source} We classify data sources into two primary types: online sampling and offline sampling. Online sampling indicates that the training data are generated through interactions using the current policy model during training, whereas offline sampling implies that training data are drawn from outputs produced by the pre-existing SFT model. Among the methods compared, RFT and DPO utilize offline sampling strategies, while Online RFT and GRPO employ online sampling techniques.

Figure \ref{fig:rl-analysis} illustrates our key finding that Online RFT achieves markedly better results than RFT across both benchmarks. More precisely, during the early training phases, Online RFT's performance is comparable to RFT, but it surpasses RFT considerably as training progresses, underscoring the efficacy of online data collection. This phenomenon can be readily explained: initially, the actor and SFT model remain highly similar, so data collected reveal only slight discrepancies. As training continues, the actor’s outputs diverge more substantially, allowing online sampling to capture richer and more informative data, thus conferring significant benefits.

\paragraph{Remarks on Gradient Coefficient} The algorithms transform the reward feedback into gradient coefficients that guide parameter updates. In our analysis, reward functions are categorized as either ‘Rule’-based or ‘Model’-based. A ‘Rule’ approach assesses responses using explicit correctness criteria, whereas the ‘Model’ approach involves a learned reward model that assigns scores to responses, trained from data labeled by rule-based evaluations. Equations \ref{eq:GC-RFT} and \ref{eq:GC-GRPO} emphasize a critical distinction: in GRPO, the gradient coefficient is dynamically scaled according to the reward value from the reward model, thereby enabling more nuanced promotion or suppression of model responses based on the degree of reward. By contrast, Online RFT does not exhibit this adaptive scaling; it offers no penalty for incorrect answers and reinforces all correct responses with equal strength.

As illustrated in Figure \ref{fig:rl-analysis}, GRPO outperforms online RFT, underscoring the advantage of modulating positive and negative gradient coefficients. Moreover, GRPO+PS achieves better results than GRPO+OS, underscoring the efficacy of fine-grained, step-specific gradient coefficients. Additionally, we investigate the impact of iterative RL by running two rounds of reinforcement learning in our experiments. According to Figure \ref{fig:iter-rl}, the introduction of iterative RL leads to a marked performance gain, particularly during the initial iteration.

\subsubsection{Why RL Works?}  
This study applies reinforcement learning to a subset of instruction tuning data, which leads to notable improvements over the baseline instruction tuning model. To further elucidate the mechanism underlying the success of reinforcement learning, we compare the Pass@K and Maj@K accuracies for both the Instruct and RL models across two benchmarks. The results presented in Figure \ref{fig:maj-pass} show that while RL substantially boosts Maj@K, there is no significant enhancement in Pass@K. This suggests that the effectiveness of RL primarily lies in strengthening the robustness of the output distribution, i.e., \textbf{the observed performance gains stem from increasing the likelihood of selecting the correct response among the TopK outputs, rather than from improvements to core reasoning capabilities.} In alignment with our findings, \citep{wang2023making} describe a \textbf{misalignment problem} for reasoning tasks in SFT models and demonstrate that preference alignment methods can improve SFT reasoning performance \citep{yuan2023rrhf,song2023preference,wang2023making}.

\subsubsection{How to Achieve More Effective RL?}
\label{sec:effec_rl}
Our experiments confirm that reinforcement learning yields strong results on mathematical reasoning tasks. Additionally, we introduce a unified framework to interpret representative training methodologies, recasting them within the context of direct or approximate RL approaches. According to Equation \ref{eq:objective}, the essential factors are the Data Source, Algorithm, and Reward Function. Below, we discuss several prospective avenues related to these components.

\paragraph{Data Source} The data source underpins every training approach. In RL scenarios, this refers specifically to the set of unlabeled questions accompanied by samples drawn from the current policy model. In our implementation, the data set consists exclusively of questions used during instruction tuning, with outputs generated via straightforward nucleus sampling. This limitation could explain why our RL pipeline primarily elevates Maj@K performance. Looking forward, we intend to extend our RL framework to incorporate out-of-distribution question prompts, and integrate \textbf{advanced sampling (decoding) strategies}, such as those based on tree-search methodologies \citep{yao2023tree}. Additionally, the role of \textbf{efficient inference techniques} \citep{xia-etal-2023-speculative,leviathan2023fast,kwon2023efficient,xia2024unlocking}, which are crucial for effective exploration of policy models, will be increasingly emphasized.

\paragraph{Algorithms} Algorithms utilize both data and reward signals to compute gradient coefficients that adjust model parameters accordingly. According to Equation \ref{eq:objective}, virtually all current approaches fundamentally \textbf{TRUST} the reward function’s signal when modulating the conditional probability assigned to a given token—either increasing or decreasing it as indicated. Yet, in highly intricate tasks, there is no guarantee that the reward feedback is consistently accurate. For instance, even meticulously curated datasets such as PRM800K \citep{lightman2023let}—annotated by trained professionals—still display an error rate of approximately 20\% in their labels\footnote{\url{https://github.com/openai/prm800k/issues/12\#issuecomment-1728491852}}. In light of these challenges, we seek to investigate reinforcement learning algorithms capable of maintaining robustness amid noisy or imperfect reward signals. Our hypothesis is that \textbf{WEAK-TO-STRONG} alignment strategies \citep{burns2023weak} could introduce transformative advancements in algorithmic learning paradigms.

\paragraph{Reward Function} The reward function delivers the essential signal that guides training. In reinforcement learning, this is typically implemented as a neural reward model. We identify three critical avenues for developing reward models: 1) \textbf{Improving the reward model's generalization capacity.} Reward models must reliably generalize to out-of-distribution prompts and sophisticated outputs, otherwise RL may simply reinforce existing distributions of LLMs without substantive improvements to their core competencies; 2) \textbf{Capturing and quantifying the reward model's uncertainty.} Proper representation of uncertainty could serve as a bridge between underpowered reward models and the weak-to-strong learning strategies; 3) \textbf{Building efficient, high-quality process reward models} that can deliver nuanced and granular feedback throughout the reasoning process \citep{lightman2023let,wang2023math}.

\section{Conclusion, Limitation, and Future Work} In this work, we introduce DeepSeekMath, an open-source model that achieves superior results over existing open-source systems on the competition-level MATH benchmark and attains performance levels close to those of proprietary models. DeepSeekMath~is built upon DeepSeek-Coder-v1.5 7B and undergoes continual training across 500B tokens, which notably includes 120B mathematics-related tokens extracted from Common Crawl. Our comprehensive ablation experiments indicate that web-sourced pages can serve as rich reservoirs for high-quality mathematical content, whereas data from arXiv appears less advantageous than anticipated. We also propose Group Relative Policy Optimization (GRPO), a variant of Proximal Policy Optimization (PPO), that enhances mathematical reasoning ability while being more memory efficient. Experimental evidence confirms that even as DeepSeekMath-Instruct 7B achieves high benchmark scores, GRPO continues to provide effective improvements. Additionally, we unify several methodologies under a common framework and highlight various avenues for advancing reinforcement learning efficacy.

Despite the strong results of DeepSeekMath~on quantitative reasoning assessments, its performance on geometry tasks and theorem-proving remains comparatively below that of closed models. Notably, during preliminary testing, the model was unable to solve problems pertaining to triangles and ellipses, possibly reflecting selection bias in its pre-training and fine-tuning datasets. Furthermore, the model’s limited scale results in reduced few-shot learning ability relative to GPT-4, which benefits more from few-shot scenarios. DeepSeekMath~shows comparable outcomes between zero-shot and few-shot settings. Moving forward, we plan to refine our data selection strategy to curate even higher-quality pre-training datasets. We also intend to investigate further opportunities for improving reinforcement learning in LLMs, as outlined in Section \ref{sec:effec_rl}.